\begingroup
\renewcommand{\arraystretch}{1.45}
\clearpage\fancyhead[L]{\small LEMNISCATIC FUNCTIONS I.}
\par\Needspace{8\baselineskip}\pdfbookmark[0]{Some New Formulas Concerning Lemniscatic Functions}{lemn-new-formulas}\begin{center}\textit{Some New Formulas Concerning Lemniscatic Functions.}\end{center}\selectlanguage{english}
Let $\displaystyle\int\frac{\mathrm{d}x}{\surd(1-x^4)}=\varphi$, or
\Needspace{10\baselineskip}
\[
\begin{aligned}\varphi&=x+\dfrac12.\dfrac15x^5+\dfrac{1.3}{2.4}.\dfrac19x^9\\[6pt]&\quad{}+\dfrac{1.3.5}{2.4.6}\dfrac1{13}x^{13}\mathord{.}\mathord{.}\end{aligned}
\]

Then one has
\Needspace{10\baselineskip}
\[
\begin{aligned}\varphi\varphi&=xx+\dfrac35.\dfrac13x^6+\dfrac{3.7}{5.9}.\dfrac15x^{10}\\[6pt]&\quad{}+\dfrac{3.7.11}{5.9.13}.\dfrac17x^{14}.\end{aligned}
\]
Let $x=\slemn\varphi$, $y=\slemn\psi$, $z=\slemn(\varphi+\psi)$

then one has
\Needspace{10\baselineskip}
\[
\begin{aligned}z&=\dfrac{x\surd(1-y^4)+y\surd(1-x^4)}{1+xxyy}\\[6pt]&=\dfrac{xx-yy}{x\surd(1-y^4)-y\surd(1-x^4)}\end{aligned}
\]
\Needspace{10\baselineskip}
\[
\begin{aligned}\surd\dfrac{1-zz}{1+zz}&=\dfrac{-2xy+\surd(1-x^4)\surd(1-y^4)}{1+xx+yy-xxyy}\\[6pt]&=\dfrac{1-xx-yy-xxyy}{2xy+\surd(1-x^4)\surd(1-y^4)}\end{aligned}
\]
\Needspace{10\baselineskip}
\[
\begin{aligned}\surd\dfrac{1-z^4}{zz}&=\dfrac{\begin{gathered}x(1+y^4)\surd(1-x^4)\\[3pt]{}-y(1+x^4)\surd(1-y^4)\end{gathered}}{(1+xxyy)(xx-yy)}\end{aligned}
\]
\Needspace{10\baselineskip}
\[
\begin{aligned}\surd(1-z^4)&=\dfrac{\begin{gathered}(1-xxyy)\surd(1-x^4)\surd(1-y^4)\\[3pt]{}-2xy(xx+yy)\end{gathered}}{(1+xxyy)^2}\end{aligned}
\]
\Needspace{10\baselineskip}
\[
\begin{aligned}\surd(1-zz)&=\dfrac{\begin{gathered}\surd(1-xx)\surd(1-yy)\\[3pt]{}-xy\surd(1+xx)\surd(1+yy)\end{gathered}}{1+xxyy}\end{aligned}
\]
\Needspace{10\baselineskip}
\[
\begin{aligned}\surd(1+zz)&=\dfrac{\begin{gathered}\surd1+xx)\surd(1+yy)\\[3pt]{}+xy\surd(1-xx)\surd(1-yy)\end{gathered}}{1+xxyy}\end{aligned}
\]

\SourceLeaf{416}
\Needspace{16\baselineskip}
The following seems to be the simplest way to expand $\slemn\varphi$ in powers of $\varphi$: one has
\Needspace{8\baselineskip}
\[
\slemn(1+i)\varphi=\dfrac{(1+i)\slemn\varphi}{\surd(1-\slemn\varphi^4)}
\]
thus, setting
\Needspace{8\baselineskip}
\[
(\slemn\varphi)^{-2}=\varphi^{-2}(1+\alpha\varphi^4+\beta\varphi^8+\gamma\varphi^{12}\mathord{.}\mathord{.}\mathord{.}\mathord{.})
\]
and consequently
\Needspace{10\baselineskip}
\[
\begin{aligned}(\slemn(1+i)\varphi)^{-2}&=-\tfrac12i\varphi^{-2}(1-4\alpha\varphi^4+16\beta\varphi^8\\[6pt]&\quad{}-64\gamma\varphi^{12}\mathord{.}\mathord{.}\mathord{.})\end{aligned}
\]
hence
\Needspace{10\baselineskip}
\[
\begin{aligned}
\slemn\varphi^2&=-2i(\slemn(1+i)\varphi)^{-2}+(\slemn\varphi)^{-2}\\[6pt]
&=5\alpha\varphi\varphi-15\beta\varphi^6+65\gamma\varphi^{10}-255\delta\varphi^{14}\mathord{.}\mathord{.}\mathord{.}
\end{aligned}
\]
one therefore has
\Needspace{10\baselineskip}
\[
\begin{aligned}(1+\alpha t+\beta t^2+\gamma t^3+\delta t^4+\mathord{.}\mathord{.}\mathord{.})\\[6pt]&\quad(5\alpha-15\beta t+65\gamma t^2-255\delta t^3\\[5pt]&\quad{}+1025\varepsilon t^4\mathord{.}\mathord{.}\mathord{.})=1\end{aligned}
\]
\Needspace{10\baselineskip}
from this
\[
\begin{aligned}5\alpha=1\\[5pt]&\alpha=\dfrac15\\[5pt]&\quad=\dfrac15\end{aligned}
\]
\Needspace{10\baselineskip}
\[
\begin{aligned}3\beta=\alpha\alpha\\[5pt]&\beta=\dfrac1{75}\\[5pt]&\quad=\dfrac1{3.25}\\[5pt]&\quad=\dfrac1{15}\alpha\end{aligned}
\]
\Needspace{10\baselineskip}
\[
\begin{aligned}13\gamma=2\alpha\beta\\[5pt]&\gamma=\dfrac2{4875}\\[5pt]&\quad=\dfrac2{3.5^3.13}\\[5pt]&\quad=\dfrac2{65}\beta\end{aligned}
\]
\Needspace{10\baselineskip}
\[
\begin{aligned}51\delta=14\alpha\gamma-3\beta\beta\\[5pt]&\delta=\dfrac1{82875}\\[5pt]&\quad=\dfrac1{3.5^3.13.17}\\[5pt]&\quad=\dfrac1{34}\gamma\end{aligned}
\]
\Needspace{10\baselineskip}
\[
\begin{aligned}205\varepsilon=50\alpha\delta-10\beta\gamma\\[5pt]&\varepsilon=\dfrac2{6215625}\\[5pt]&\quad=\dfrac2{9.5^5.13.17}\\[5pt]&\quad=\dfrac2{75}\delta\end{aligned}
\]
\Needspace{10\baselineskip}
\[
\begin{aligned}819\zeta=206\alpha\varepsilon-54\beta\delta+13\gamma\gamma\\[5pt]&\zeta=\dfrac2{242409375}\\[5pt]&\quad=\dfrac2{3^3.5^5.13^2.17}\\[5pt]&\quad=\dfrac1{39}\varepsilon\end{aligned}
\]
\Needspace{10\baselineskip}
\[
\begin{aligned}3277\eta=818\alpha\zeta-202\beta\varepsilon+38\gamma\delta\\[5pt]&\eta=\dfrac4{19527421875}\\[5pt]&\quad=\dfrac4{3.5^7.13^2.17.29}\\[5pt]&\quad=\dfrac{18}{725}\zeta\end{aligned}
\]
\Needspace{10\baselineskip}
\[
\begin{aligned}13107\theta=3278\alpha\eta-822\beta\zeta+218\gamma\varepsilon-51\delta\delta,\\[5pt]&\theta=\dfrac{223}{44815433203125}\\[5pt]&\quad=\dfrac{223}{3^4.5^8.13^2.17^2.29}\\[5pt]&\quad=\dfrac{223}{9180}\eta\end{aligned}
\]
The limit of the ratio of two successive terms is
\Needspace{6\baselineskip}
\[
(2{,}6220\mathord{.}\mathord{.})^4:1=47{,}27:1
\]

\SourceLeaf{417}

To a very close approximation,
\Needspace{8\baselineskip}
\[
\zeta=\dfrac{92}{(2\varpi)^{24}},\qquad
\eta=\dfrac{108}{(2\varpi)^{28}},\qquad
\theta=\dfrac{124}{(2\varpi)^{32}}
\]
\Needspace{22\baselineskip}
\[
\begin{aligned}
\log P&=\log\varphi-\dfrac1{12}\alpha\varphi^4-\dfrac1{56}\beta\varphi^8
\\[6pt]&\quad{}-\dfrac1{132}\gamma\varphi^{12}-\dfrac1{240}\delta\varphi^{16}
\\[6pt]&\quad{}-\dfrac1{380}\varepsilon\varphi^{20}-\dfrac1{552}\zeta\varphi^{24}\mathord{.}\mathord{.}\mathord{.}\\[6pt]
&=\log\varphi-\dfrac1{60}\varphi^4-\dfrac1{4200}\varphi^8
\\[6pt]&\quad{}-\dfrac1{321750}\varphi^{12}-\dfrac1{19890000}\varphi^{16}
-\dfrac1{1180968750}\varphi^{20}\\[6pt]
&\quad-\dfrac{22}{1756255921875}\varphi^{24}-\mathord{.}\mathord{.}
\end{aligned}
\]
\Needspace{22\baselineskip}
\[
\begin{aligned}
\log\operatorname{sin\,lemn}\varphi
&=\log\varphi-\dfrac36\alpha\varphi^4+\dfrac7{28}\beta\varphi^8
\\[6pt]&\quad{}-\dfrac{33}{66}\gamma\varphi^{12}+\dfrac{127}{120}\delta\varphi^{16}
\\[6pt]&\quad{}-\dfrac{513}{190}\varepsilon\varphi^{20}+\dfrac{2047}{276}\zeta\varphi^{24}\mathord{.}\mathord{.}\\[6pt]
&=\log\varphi-\dfrac1{10}\varphi^4+\dfrac1{300}\varphi^8
\\[6pt]&\quad{}-\dfrac1{4875}\varphi^{12}+\dfrac{127}{9945000}\varphi^{16}
-\dfrac3{3453125}\varphi^{20}\\[6pt]
&\quad+\dfrac{89}{1454456250}\varphi^{24}-\mathord{.}\mathord{.}
\end{aligned}
\]
\Needspace{14\baselineskip}
The coefficients $\alpha$, $\beta$, $\gamma$, and so on can also be determined by the following equation
\[
\begin{aligned}1+\alpha t+\beta tt+\gamma t^3+\mathord{.}\mathord{.}\mathord{.}
\\[6pt]&=\left(\begin{gathered}1+\dfrac3{2.3}\alpha t+\dfrac7{8.7}\beta tt
\\[4pt]{}+\dfrac{33}{32.11}\gamma t^3+\dfrac{127}{128.15}\delta t^4\mathord{.}\mathord{.}\mathord{.}\end{gathered}\right)^2\end{aligned}
\]
The most convenient way to expand $\log\operatorname{cos\,lemn}\varphi$ in a series is as follows. One has
\Needspace{10\baselineskip}
\[
\begin{aligned}\dfrac{\mathrm d\log\operatorname{cos\,lemn}\varphi}{\mathrm d\varphi}
&=-\dfrac{2\mathrm d\varphi'}{\mathrm d\log\operatorname{sin\,lemn}\varphi}
\\[7pt]&=-(1-i)\operatorname{sin\,lemn}(1+i)\varphi\end{aligned}
\]
this gives
\Needspace{18\baselineskip}
\[
\begin{aligned}
\log\operatorname{cos\,lemn}\varphi
&=-\varphi\varphi-\dfrac2{15}\varphi^6-\dfrac2{75}\varphi^{10}
\\[6pt]&\quad{}-\dfrac{44}{6825}\varphi^{14}-\dfrac{422}{248625}\varphi^{18}
\\[6pt]&\quad{}-\dfrac{6428}{13673375}\varphi^{22}\\[6pt]
&\quad-\dfrac{6044}{44890625}\varphi^{26}
-\dfrac{20824792}{527240390625}\varphi^{30}-\mathord{.}\mathord{.}
\end{aligned}
\]
\Needspace{8\baselineskip}
\[
P\dfrac{\mathrm d^4P}{\mathrm d\varphi^4}
-4\dfrac{\mathrm dP}{\mathrm d\varphi}\, .\,\dfrac{\mathrm d^3P}{\mathrm d\varphi^3}
+3\left(\dfrac{\mathrm d\mathrm dP}{\mathrm d\varphi^2}\right)^2=2PP
\]
\Needspace{24\baselineskip}
\[
\begin{aligned}
P\varphi&=\varphi-\dfrac2{1\mathord{.}\mathord{.}\mathord{.}\mathord{.}5}\varphi^5
-\dfrac{36}{1\mathord{.}\mathord{.}\mathord{.}\mathord{.}9}\varphi^9
\\[6pt]&\quad{}+\dfrac{552}{1\mathord{.}\mathord{.}\mathord{.}\mathord{.}13}\varphi^{13}
+\dfrac{5136}{1\mathord{.}\mathord{.}\mathord{.}\mathord{.}17}\varphi^{17}
\\[6pt]&\quad{}+\dfrac{5146848}{1\mathord{.}\mathord{.}\mathord{.}\mathord{.}21}\varphi^{21}-\mathord{.}\mathord{.}\\[6pt]
&=\varphi-\dfrac1{4.3.5}\varphi^5-\dfrac1{32.9.5.7}\varphi^9
\\[6pt]&\quad{}+\dfrac{23}{128.81.25.7.11.13}\varphi^{13}
\\[6pt]&\quad{}+\dfrac{107}{2048.243.125.49.11.13.17}\varphi^{17}\\[6pt]
&\quad+\dfrac{23.37}{8192.729.625.49.11.13.17.19}\varphi^{21}-\mathord{.}\mathord{.}
\end{aligned}
\]
\SourceLeaf{418}
\Needspace{12\baselineskip}
\[
\begin{aligned}
P\psi\varpi&=+2.6220575\,\psi\\
&\phantom{{}={}}-2.0656648\,\psi^5\\
&\phantom{{}={}}-0.5811918\,\psi^9\\
&\phantom{{}={}}+0.0245475\,\psi^{13}\\
&\phantom{{}={}}+0.0001890\,\psi^{17}\\
&\phantom{{}={}}+0.0000620\,\psi^{21}
\end{aligned}
\]
\Needspace{8\baselineskip}
\[
P\tfrac12\varpi=+1.2453446
={}^{4}\!\surd\tfrac12\, .\,e^{\dfrac18\pi}
\]
\Needspace{8\baselineskip}
\[
\dfrac{\mathrm d\mathrm d\log Q}{\mathrm d\varphi^2}=\dfrac{PP}{QQ}
\]
\Needspace{18\baselineskip}
Halving can be done very conveniently as follows
\Needspace{16\baselineskip}
\[
\begin{aligned}
\operatorname{cos\,lemn}\varphi&=\operatorname{tang}u\\[8pt]
\operatorname{sin\,lemn}\varphi&=\surd\cos2u\\[8pt]
\operatorname{cos\,lemn}\varphi&=\surd\operatorname{tang}\tfrac12(45^0+u)\\[8pt]
\operatorname{sin\,lemn}\varphi&=\surd\operatorname{tang}\tfrac12(45^0-u)
\end{aligned}
\]
\Needspace{18\baselineskip}
\[
\begin{aligned}
\operatorname{sin\,lemn}(\varphi+i\psi)&=r(\cos v+i\sin v)\\[8pt]
\operatorname{sin\,lemn}\varphi^2&=\operatorname{tang}\Phi\\[8pt]
\operatorname{sin\,lemn}\psi^2&=\operatorname{tang}\Psi\\[8pt]
r&=\surd\operatorname{tang}(\Phi+\Psi)\\[8pt]
\operatorname{tang}v&=\surd\dfrac{\operatorname{tang}2\Psi}{\operatorname{tang}2\Phi}
\end{aligned}
\]
\Needspace{16\baselineskip}
\[
\begin{aligned}
\operatorname{sin\,lemn}(\varphi+\psi)
&=\dfrac{\surd(\cos2\Phi\sin2\Psi)+\surd(\sin2\Phi\cos2\Psi)}{\cos(\Phi-\Psi)\, .\,\surd2}\\[8pt]
&=\dfrac{\sin(\Psi-\Phi)\, .\,\surd2}{\surd(\cos2\Phi\sin2\Psi)-\surd(\sin2\Phi\cos2\Psi)}
\end{aligned}
\]
\Needspace{16\baselineskip}
Let
\Needspace{12\baselineskip}
\[
\operatorname{sin\,lemn}A=x,\qquad
\operatorname{sin\,lemn}B=y,\qquad
\operatorname{sin\,lemn}\dfrac{A+B}{1+i}=z
\]
then
\Needspace{10\baselineskip}
\[
\begin{aligned}(1+i)z&=\dfrac{\begin{gathered}\surd(1+xx)(1-yy)\\[3pt]{}-\surd(1-xx)(1+yy)\end{gathered}}{x-y}\end{aligned}
\]
\Needspace{10\baselineskip}
\[
\begin{aligned}\dfrac{1-i}{z}&=\dfrac{\begin{gathered}\surd1+xx)(1-yy)\\[3pt]{}+\surd(1-xx)(1+yy)\end{gathered}}{x+y}\end{aligned}
\]
\SourceLeaf{419}
\Needspace{10\baselineskip}
\[
\dfrac{\operatorname{sin\,lemn}\dfrac{A-B}{1+i}}
{\operatorname{sin\,lemn}\dfrac{A+B}{1+i}}
=\dfrac{\operatorname{sin\,lemn}A-\operatorname{sin\,lemn}B}
{\operatorname{sin\,lemn}A+\operatorname{sin\,lemn}B}
\]
\Needspace{8\baselineskip}
\[
pp=QQ-PP,\qquad qq=QQ+PP
\]
\Needspace{10\baselineskip}
\[
\begin{aligned}P(a-b)\, .\,Q(a+b)&=Pa\, .\,Qa\, .\,\surd(Qb^4-Pb^4)\\[5pt]&\quad{}-Pb\, .\,Qb\, .\,\surd(Qa^4-Pa^4)\end{aligned}
\]
\Needspace{8\baselineskip}
\[
\begin{aligned}P(a-b)\, .\,P(a+b)&=Pa^2\, .\,Qb^2-Qa^2\, .\,Pb^2\end{aligned}
\]
\Needspace{8\baselineskip}
\[
\begin{aligned}Q(a-b)\, .\,Q(a+b)&=Pa^2\, .\,Pb^2+Qa^2\, .\,Qb^2\end{aligned}
\]
\Needspace{10\baselineskip}
\[
\begin{aligned}
q(a-b)\, .\,q(a+b)&=pa^2\, .\,Pb^2+qa^2\, .\,Qb^2\\
&=pb^2\, .\,Pa^2+qb^2\, .\,Qa^2
\end{aligned}
\]
\Needspace{10\baselineskip}
\[
\begin{aligned}q(a-b)\, .\,Q(a+b)&=qa\, .\,qb\, .\,Qa\, .\,Qb\\[5pt]&\quad{}-pa\, .\,pb:Pa\, .\,Pb\end{aligned}
\]
\Needspace{10\baselineskip}
\[
\begin{aligned}q(a-b)\, .\,P(a+b)&=qa\, .\,pb\, .\,Pa\, .\,Qb\\[5pt]&\quad{}+pa\, .\,qb\, .\,Qa\, .\,Pb\end{aligned}
\]
\Needspace{14\baselineskip}
\[
\begin{aligned}
P\varphi&=P\\[4pt]
Q\varphi&=Q\\[4pt]
p\varphi&=\surd(QQ-PP)\\[4pt]
q\varphi&=\surd(QQ+PP)
\end{aligned}
\]
\Needspace{14\baselineskip}
\[
\begin{aligned}
P2\varphi&=2PQ\surd(Q^4-P^4)\\[4pt]
Q2\varphi&=Q^4+P^4\\[4pt]
p2\varphi&=Q^4-2QQPP-P^4\\[4pt]
q2\varphi&=Q^4+2QQPP+P^4
\end{aligned}
\]
\Needspace{8\baselineskip}
\[
\begin{aligned}P3\varphi&=3Q^8P-6Q^4P^5-P^9\end{aligned}
\]
\Needspace{8\baselineskip}
\[
\begin{aligned}Q3\varphi&=Q^9+6Q^5P^4-3QP^8\end{aligned}
\]
\Needspace{10\baselineskip}
\[
\begin{aligned}p3\varphi&=\surd(QQ-PP)\, .\,(Q^8-4Q^6PP\\[5pt]&\quad{}-6Q^4P^4-4QQP^6+P^8)\end{aligned}
\]
\Needspace{10\baselineskip}
\[
\begin{aligned}q3\varphi&=\surd(QQ+PP)\, .\,(Q^4+4Q^6PP\\[5pt]&\quad{}-6Q^4P^4+4QQP^6+P^8)\end{aligned}
\]
\Needspace{12\baselineskip}
\[
\begin{aligned}P4\varphi&=4PQ\surd(Q^4-P^4)\\[5pt]&\quad{}\, .\,(Q^{12}-5Q^8P^4-5Q^4P^8+P^{12})\end{aligned}
\]
\Needspace{12\baselineskip}
\[
\begin{aligned}Q4\varphi&=Q^{16}+20Q^{12}P^4-26Q^8P^8\\[5pt]&\quad{}+20Q^4P^{12}+P^{16}\end{aligned}
\]
\Needspace{14\baselineskip}
\[
\begin{aligned}p4\varphi&=Q^{16}-8Q^{14}PP-12Q^{12}P^4\\[5pt]&\quad{}-8Q^{10}P^6+38Q^8P^8+8Q^6P^{10}\\[5pt]&\quad{}-12Q^4P^{12}+8QQP^4+P^{16}\end{aligned}
\]
\Needspace{14\baselineskip}
\[
\begin{aligned}q4\varphi&=Q^{16}+8Q^{14}PP-12Q^{12}P^4\\[5pt]&\quad{}+8Q^{10}P^6+38Q^8P^8-8Q^6P^{10}\\[5pt]&\quad{}-12Q^4P^{12}-8QQP^4+P^{16}\end{aligned}
\]


\SourceLeaf{420}
\Needspace{30\baselineskip}
\[
\begin{aligned}
\dfrac{Qn\varphi}{Q\varphi^{nn}}
&=1+\dfrac1{12}(n^4-nn)\left(\dfrac{P\varphi}{Q\varphi}\right)^4
\\[5pt]&\quad{}-\dfrac1{10080}(n^8+70n^6\\[5pt]&\quad{}-371n^4+300nn)\left(\dfrac{P\varphi}{Q\varphi}\right)^8\\[7pt]
&\quad+\dfrac1{19958400}(17n^{12}+165n^{10}\\[5pt]&\quad{}+4191n^8-106865n^6+426492n^4\\[7pt]
&\quad-324000nn)\, .\,\left(\dfrac{P\varphi}{Q\varphi}\right)^{12}\\[7pt]
&=1+\dfrac1{12}nn(n-1)\left(\dfrac{P\varphi}{Q\varphi}\right)^4
\\[5pt]&\quad{}-\dfrac1{10080}nn(nn-1)(nn-4)(nn+75)\left(\dfrac{P\varphi}{Q\varphi}\right)^8\\[7pt]
&\quad+\dfrac1{19958400}nn(nn-1)(nn-4)(nn-9)\\[5pt]&\quad{}(17n^4+403nn+9000)\, .\,\left(\dfrac{P\varphi}{Q\varphi}\right)^{12}
\end{aligned}
\]
\Needspace{14\baselineskip}
\[
\begin{aligned}
Pi\varphi&=iP\varphi\\[4pt]
Qi\varphi&=Q\varphi\\[4pt]
pi\varphi&=q\varphi\\[4pt]
qi\varphi&=p\varphi
\end{aligned}
\]
\Needspace{14\baselineskip}
\[
\begin{aligned}
P(1+i)\varphi&=(1+i)PQ\\[4pt]
Q(1+i)\varphi&=\surd(Q^4-P^4)=pq\\[4pt]
p(1+i)\varphi&=QQ-iPP\\[4pt]
q(1+i)\varphi&=QQ+iPP
\end{aligned}
\]
\Needspace{8\baselineskip}
\[
\begin{aligned}P(2+i)\varphi&=(2+i)PQ^4-iP^5\end{aligned}
\]
\Needspace{8\baselineskip}
\[
\begin{aligned}Q(2+i)\varphi&=Q^5-(1-2i)QP^4\end{aligned}
\]
\Needspace{10\baselineskip}
\[
\begin{aligned}p(2+i)\varphi&=\surd(QQ+PP)\\[5pt]&\quad{}\, .\,\{Q^4-(2+2i)QQPP+P^4\}\end{aligned}
\]
\Needspace{10\baselineskip}
\[
\begin{aligned}q(2+i)\varphi&=\surd(QQ-PP)\\[5pt]&\quad{}\, .\,\{Q^4+(2+2i)QQPP+P^4\}\end{aligned}
\]
\Needspace{10\baselineskip}
\[
\begin{aligned}P(3+i)\varphi&=(3+i)PQ^3-(2+6i)P^5Q^5\\[5pt]&\quad{}+(3+i)P^9Q\end{aligned}
\]
\Needspace{10\baselineskip}
\[
\begin{aligned}Q(3+i)\varphi&=\surd(Q^4-P^4)\\[5pt]&\quad{}\, .\,\{Q^8+(2+8i)P^4Q^4+P^8\}\end{aligned}
\]
\Needspace{14\baselineskip}
\[
\begin{aligned}p(3+i)\varphi&=Q^{10}-(4-3i)Q^8PP\\[5pt]&\quad{}-(2-4i)Q^6P^4+(4+2i)Q^4P^6\\[5pt]&\quad{}-(3+4i)QQP^8-iP^{10}\end{aligned}
\]
\Needspace{14\baselineskip}
\[
\begin{aligned}q(3+i)\varphi&=Q^{10}+(4-3i)Q^8PP\\[5pt]&\quad{}-(2-4i)Q^6P^4-(4+2i)Q^4P^6\\[5pt]&\quad{}-(3+4i)QQP^8+iP^{10}\end{aligned}
\]
\Needspace{8\baselineskip}
\[
\begin{aligned}P(1+2i)\varphi&=(1+2i)PQ^4-P^5\end{aligned}
\]
\Needspace{8\baselineskip}
\[
\begin{aligned}Q(1+2i)\varphi&=Q^5-(1+2i)QP^4\end{aligned}
\]
\Needspace{10\baselineskip}
\[
\begin{aligned}P(1+3i)\varphi&=(1+3i)PQ^9-(6+2i)P^5Q^5\\[5pt]&\quad{}+(1+3i)P^9Q\end{aligned}
\]
\Needspace{6\baselineskip}
\[
\begin{aligned}Q(1+3i)\varphi&=Q^{10}\mathord{.}\mathord{.}\end{aligned}
\]
\Needspace{14\baselineskip}
\[
\begin{aligned}P(1+4i)\varphi&=(1+4i)PQ^{16}-(20+12i)P^5Q^{12}\\[5pt]&\quad{}-(10-28i)P^9Q^8\\[5pt]&\quad{}+(12-20i)P^{13}Q^4+P^{17}\end{aligned}
\]
\Needspace{6\baselineskip}
\[
\begin{aligned}Q(1+4i)\varphi&=\end{aligned}
\]

\SourceLeaf{421}
\Needspace{14\baselineskip}
\[
\begin{aligned}
P(1+ni)\varphi
 &= (1+in)PQ^{nn}
 \\[5pt]&\quad{}-\left(\dfrac{nn.nn-1}{12}+\dfrac{n.nn-1.nn-4}{60}i\right)P^5Q^{nn-4}+{}\\[7pt]
 &= (1+in)PQ^{nn}\\[5pt]&\quad{}-\dfrac{n.nn-1.1+in.n-4i}{60}P^5Q^{nn-4}
\end{aligned}
\]
Among the numbers $y=x,\ 2x+1,\ 2x+i,\ 2x-1,\ 2x-i$ there is always at least one (or three or all) that makes it possible to solve the congruence $1-y^4\equiv zz$ for any modulus.
\Needspace{28\baselineskip}
\[
\slemn X=x,\qquad \slemn Y=y,\qquad \slemn Z=z
\]
\Needspace{22\baselineskip}
\[
\begin{gathered}\slemn(X+Y+Z)=\\[7pt]
\dfrac{\begin{aligned}
 &-2xyz(xx+yy+zz-xxyyzz)\\[4pt]&\quad{}+x\surd(1-y^4).\surd(1-z^4)\\[3pt]&\quad{}.(1-yyzz+xxzz+xxyy)\\
 &+y\surd(1-x^4).\surd(1-z^4)\\[3pt]&\quad{}.(1+yyzz-xxzz+xxyy)\\
 &+z\surd(1-x^4).\surd(1-y^4)\\[3pt]&\quad{}.(1+yyzz+xxzz-xxyy)
\end{aligned}}{\begin{gathered}1+2yyzz+2xxzz+2xxyy+y^4z^4\\[4pt]{}+x^4z^4+x^4y^4-2x^4yyzz\\[4pt]{}-2xxy^4zz-2xxyyz^4\end{gathered}}\end{gathered}
\]

\Needspace{18\baselineskip}
If
\Needspace{14\baselineskip}
\[
\begin{array}{rll}
\varphi+\varphi'+\varphi''=0,\\[5pt]&\slemn\varphi=x&\clemn\varphi=y\\[4pt]
&\slemn\varphi'=x'&\clemn\varphi'=y'\\[4pt]
&\slemn\varphi''=x''&\clemn\varphi'=y''
\end{array}
\]
then
\Needspace{10\baselineskip}
\[
\begin{aligned}x(y-y'y'')&=x'(y'-yy'')=x''(y''-yy'),\\[7pt]&\qquad
y''=\dfrac{xy-x'y'}{xy'-x'y}\end{aligned}
\]
\Needspace{24\baselineskip}
\begin{center}
\small
\begin{tabular}{r|r|l}
\multicolumn{1}{c}{$Q$}&\multicolumn{1}{c}{}&\multicolumn{1}{c}{}\\
$-1+2i$&$5$&$1+(-1+2i)x^4$\\[9pt]
$-3$&$9$&$1+6x^4-3x^8$\\[9pt]
$+3+2i$&$13$&$\begin{aligned}[t]1+(-11+10i)x^4\\[3pt]{}+(7-4i)x^8+(3+2i)x^{12}\end{aligned}$\\[9pt]
$+1+4i$&$17$&$\begin{aligned}[t]1+(+12-20i)x^4\\[3pt]{}+(-10+28i)x^8-(20+12i)x^{12}\\[3pt]{}+(1+4i)x^{16}\end{aligned}$\\[9pt]
$5$&$25$&$\begin{aligned}[t]1+50x^4-125x^8\\[3pt]{}+300x^{12}-105x^{16}\\[3pt]{}-62x^{20}+5x^{24}\end{aligned}$
\end{tabular}
\end{center}

\Needspace{14\baselineskip}
for \hfill $x^4=+1$\quad and\quad $x^4=-1$\hfill\mbox{}

these functions become respectively equal to the squares and cubes
\Needspace{8\baselineskip}
\[
\begin{array}{rccccc}
\text{of}&1-i&-2&-2-2i&-4i&+8\\
\text{for }M=&5&9&13&17&25
\end{array}
\]
\Needspace{18\baselineskip}
\[
\begin{aligned}
Q(a+bi)\varphi={}&Q^{aa+bb}+\dfrac1{12}\{(a+bi)^4-(aa+bb)\}Q^{aa+bb-4}P^4\\[7pt]
&-\dfrac1{10080}\{(a+bi)^8+70(a+bi)^4(aa+bb)\\[5pt]&\quad{}-35(aa+bb)^2-336(a+bi)^4\\[5pt]&\quad{}+300(aa+bb)\}.\\[7pt]
&\quad.Q^{aa+bb-8}P^8\mathord{.}\mathord{.}\mathord{.}
\end{aligned}
\]
\par\begin{center}\rule{.22\linewidth}{.3pt}\end{center}
\clearpage
\endgroup
